\documentclass[11pt, a4paper]{article}
\usepackage[margin=1in, top=1.2in, bottom=1.2in]{geometry}
% ---- Font & Micro-Typography ----
\usepackage[T1]{fontenc}
\usepackage{amsmath}
\usepackage{lmodern}
\usepackage{microtype}
% ---- Formatting & Aesthetics ----
\usepackage{graphicx}
% ---- Bibliography ----
\usepackage[backend=biber, style=numeric-comp, sorting=nyt, giveninits=true]{biblatex}


% ====================================================================================
\begin{document}

\begin{center}
  \LARGE\textbf{Daily Rant Fab}\\
  Tadouanla Guetchuin Billy\\
  28177
\end{center}


Today I was introduced to digital manufacturing. This is the process of creating or
manufacturing physical products using software tools and different types of machine. Fusion 360
is a \textbf{Computer Aided Design tool} with an integrated \textbf{Computer- Aided
  Manufacturing} (CAM). CAD is the use of computer software to create modify and analyse and
optimize 2D and 3D models on physical objects. CAM is the use of software and computer
controlled machinery to automate the manufacturing process.

\textbf{Design Process} relies on parametric feature-based modelling, typically
beginning with sketches on 2-dimensional planes and utilizing operations such as
extrusions, revolutions, sweeps, and more to transition from 2D profiles into complex
3D geometries.

We also had the opportunity to see one of the machines at work namely the Laser cutter. The
way it works based on my understanding is through the use of a high hit source (the beam of
light or laser) which cuts through the wood based on the design given to it.
% main section for the daily rant here 
%------ First Picture -------------
\begin{figure}[h!]
  \center
  \includegraphics[scale=0.30]{./figures/Fab_1.jpg}
\end{figure}


Below is the formula to calculate the surface area of a Cube (6 sides).
\begin{align}
  \mathbf{A = 6s^2}
\end{align}


% \newpage

%------ Second Picture -------------
% \begin{figure}[h!]
%   \center
%   \includegraphics[scale=0.5]{./figures/Fab_2.jpg}
% \end{figure}

\end{document}
